\chapter{Espaces affines}\label{ch:b2:affine}

Les espaces vectoriels possèdent un point privilégié --- l'origine ---
dont la géométrie ne veut pas. Un \emph{\omterm{def:b2:affine:def}{espace affine}} est un espace
vectoriel qui a oublié son origine~: points et vecteurs deviennent des
espèces différentes, reliées par translation. Ce court chapitre
construit le dictionnaire (points, \omterm{def:b2:affine:barycenter}{barycentres}, sous-espaces et
\omterm{def:b2:affine:subspace}{applications affines}), le point de vue \omterm{def:b2:affine:subspace}{affine} sur la convexité, et les
outils de classification utilisés dans les chapitres de géométrie à
venir.

\section{Points et vecteurs}

\begin{definition}\label{def:b2:affine:def}
Un \emph{espace affine}\index{espace affine} dirigé par un espace
vectoriel réel $E$ est un ensemble non vide $\mathcal{E}$ muni d'une
application $(A, B) \mapsto \vect{AB} \in E$ vérifiant
\[
\vect{AB} + \vect{BC} = \vect{AC}
\quad \text{(Chasles)},
\qquad
\text{pour chaque } A,\ B \mapsto \vect{AB} \text{ est une bijection }
\mathcal{E} \to E .
\]
On note $B = A + u$ l'unique point tel que $\vect{AB} = u$. La
\emph{dimension} de $\mathcal{E}$ est $\dim E$. Tout espace vectoriel
est un espace \omterm{def:b2:affine:subspace}{affine} sur lui-même ($\vect{AB} = B - A$)~; tout choix
d'origine $O \in \mathcal{E}$ identifie $\mathcal{E}$ à $E$ via $M
\mapsto \vect{OM}$.
\end{definition}

\begin{example}[Un espace affine sans origine
naturelle]\label{ex:b2:affine:planeexample}
Le plan solution $\mathcal E = \{(x, y, z) \in \R^3 : x + y
+ z = 1\}$ n'est pas un sous-espace vectoriel ($0 \notin \mathcal E$),
mais c'est un \omterm{def:b2:affine:def}{espace affine} dirigé par $E = \{x + y + z =
0\}$~: pour $A, B \in \mathcal E$ la différence $\vect{AB} = B
- A$ tombe dans $E$ (les sommes s'annulent), Chasles est hérité
de $\R^3$, et $B \mapsto \vect{AB}$ est bijective sur $E$.
Aucun point de $\mathcal E$ n'est distingué --- tout choix
d'«~origine~» $O \in \mathcal E$ convient aussi bien, et toutes les
identifications $M \mapsto \vect{OM}$ diffèrent par des translations.
C'est la situation typique~: les ensembles solutions de problèmes
linéaires inhomogènes (systèmes linéaires, équations \omterm{def:b2:diffcalc:differential}{différentielles}
linéaires au \cref{ch:b2:diffeq}) sont \omterm{def:b2:affine:subspace}{affines}, jamais linéaires,
et le slogan «~solution particulière plus noyau~» est exactement
l'énoncé $\mathcal F = A + F$ de la définition suivante.
\end{example}

\begin{definition}[Barycentre]\label{def:b2:affine:barycenter}
Soit $(A_i, \lambda_i)_{i \leq k}$ des points pondérés avec
$\sum\lambda_i \neq 0$. Le \emph{barycentre}\index{barycentre} $G
= \operatorname{bar}\bigl((A_i, \lambda_i)\bigr)$ est l'unique
point tel que
\[
\sum_i \lambda_i\, \vect{GA_i} = 0
\qquad\text{de manière équivalente}\qquad
\vect{OG} = \frac{1}{\sum\lambda_i}\sum_i \lambda_i\,\vect{OA_i}
\quad (\text{pour tout } O).
\]
Les barycentres sont \emph{associatifs} (des sous-groupes de points
peuvent être remplacés par leur barycentre partiel affecté du poids
total) et invariants par changement d'échelle de tous les poids.
\end{definition}

\begin{proof}[Démonstration de l'existence et des formules]
Fixons $O$ et posons $s = \sum_i\lambda_i \neq 0$. Par Chasles,
\[
\sum_i\lambda_i\,\vect{GA_i} = 0
\iff \sum_i\lambda_i\bigl(\vect{GO} + \vect{OA_i}\bigr) = 0
\iff s\,\vect{OG} = \sum_i\lambda_i\,\vect{OA_i},
\]
ce qui détermine $G = O + \frac1s\sum_i\lambda_i\vect{OA_i}$
de manière unique. \emph{Indépendance en $O$~:} pour une autre origine
$O'$,
\[
\frac1s\sum_i\lambda_i\,\vect{O'A_i}
= \frac1s\sum_i\lambda_i\bigl(\vect{O'O} + \vect{OA_i}\bigr)
= \vect{O'O} + \frac1s\sum_i\lambda_i\,\vect{OA_i}
= \vect{O'G} :
\]
le même point $G$. \emph{Associativité~:} partitionnons l'ensemble
d'indices en $I \sqcup J$ avec $s_I = \sum_{i\in I}\lambda_i \neq 0$,
et soit $G_I$ le \omterm{def:b2:affine:barycenter}{barycentre} de $(A_i, \lambda_i)_{i\in
I}$, de sorte que $\sum_{i\in I}\lambda_i\vect{OA_i} =
s_I\,\vect{OG_I}$. Alors
\[
s\,\vect{OG}
= \sum_{i\in I}\lambda_i\vect{OA_i} +
\sum_{j\in J}\lambda_j\vect{OA_j}
= s_I\,\vect{OG_I} + \sum_{j\in J}\lambda_j\vect{OA_j} :
\]
$G$ est le \omterm{def:b2:affine:barycenter}{barycentre} de $(G_I, s_I)$ avec $(A_j,
\lambda_j)_{j\in J}$, comme annoncé. \emph{Changement d'échelle~:}
remplacer chaque $\lambda_i$ par $t\lambda_i$ ($t \neq 0$) multiplie
$s$ et la somme pondérée par $t$, laissant $\vect{OG}$ inchangé.
\end{proof}

\begin{remark}[Pièges classiques]
Deux pièges entourent la définition. D'abord, si les poids sont de
somme \emph{nulle}, il n'y a pas de \omterm{def:b2:affine:barycenter}{barycentre}~: l'application $O
\mapsto \sum\lambda_i\vect{OA_i}$ est alors indépendante de $O$ et
définit un \emph{vecteur}, non un point --- par exemple
$(A, -1;\ B, 1)$ encode $\vect{AB}$. Garder trace de celui des deux
objets qu'un calcul produit est la moitié de l'hygiène
barycentrique. Ensuite, les poids n'ont de sens qu'à un facteur
commun non nul près~; des formules comme «~les coordonnées
de $G$ sont $\lambda_1, \dots, \lambda_k$~» présupposent une
normalisation (en général $\sum\lambda_i = 1$), et oublier
de normaliser est la source classique de mauvais rapports sur une
figure.
\end{remark}

\begin{definition}[Sous-espaces affines~; applications affines]\label{def:b2:affine:subspace}
Un \emph{sous-espace affine} est un ensemble $\mathcal{F} = A + F = \{A + u :
u \in F\}$ où $F$ est un sous-espace vectoriel (sa \emph{direction})~;
de manière équivalente, un ensemble non vide stable par \omterm{def:b2:affine:barycenter}{barycentres}.
Les sous-espaces affines de $\R^n$ sont exactement les ensembles
solutions des systèmes linéaires
$MX = B$ (année 1~: solution particulière plus noyau). Une application $f \colon
\mathcal{E} \to \mathcal{E}'$ est \emph{affine}\index{application affine}
lorsqu'elle conserve les \omterm{def:b2:affine:barycenter}{barycentres} --- de manière équivalente lorsque
\[
f(A + u) = f(A) + \varphi(u)
\]
pour une (unique) application linéaire $\varphi = \vec f$, la
\emph{partie linéaire}. Applications affines de $\R^n$~: $X \mapsto MX + C$.
Les composées sont affines, de parties linéaires composées~; $f$ est
bijective si et seulement si $\vec f$ l'est.
\end{definition}

\begin{proof}[Démonstration de l'équivalence pour les applications]
Si $f(A + u) = f(A) + \varphi(u)$~: pour un \omterm{def:b2:affine:barycenter}{barycentre} $G$ de $(A_i,
\lambda_i)$, en développant chaque point à partir de $A$, $f(G) = f(A) +
\varphi(\vect{AG})$ et $\varphi(\vect{AG}) =
\frac{\sum\lambda_i\varphi(\vect{AA_i})}{\sum\lambda_i}$~: $f(G)$ est
le \omterm{def:b2:affine:barycenter}{barycentre} des images. Réciproquement, fixons $A$ et posons
$\varphi(u) = \vect{f(A)\,f(A + u)}$. \emph{Homogénéité~:} $A + tu =
\operatorname{bar}\bigl(A, 1-t;\ A + u, t\bigr)$ pour tout réel
$t$, donc la conservation des \omterm{def:b2:affine:barycenter}{barycentres} (avec des poids réels
quelconques, comme supposé) donne directement $\varphi(tu) =
t\,\varphi(u)$. \emph{Additivité~:} $A + u + v = \operatorname{bar}\bigl(A + 2u,
\tfrac12;\ A + 2v, \tfrac12\bigr)$, donc $\varphi(u + v) =
\tfrac12\varphi(2u) + \tfrac12\varphi(2v) = \varphi(u) +
\varphi(v)$, en utilisant l'homogénéité. Donc $\varphi$ est linéaire.
\end{proof}

\begin{remark}
La démonstration a utilisé des \omterm{def:b2:affine:barycenter}{barycentres} à poids \emph{réels
quelconques}~: l'étape d'homogénéité prend $t$ hors de $\intcc01$.
Si l'on suppose seulement qu'une application conserve les \omterm{def:b2:affine:barycenter}{barycentres} à
poids positifs --- de manière équivalente, les milieux et les segments
--- la linéarité de l'application vectorielle n'est plus gratuite~:
on n'obtient que la $\Q$-linéarité, et une hypothèse de continuité est
nécessaire pour conclure, exactement comme à
l'\cref{exo:b2:affine:5}. Distinguer «~conserve tous les
\omterm{def:b2:affine:barycenter}{barycentres}~» de «~conserve les combinaisons \omterm{def:b2:affine:convex}{convexes}~» est une
subtilité petite mais réelle du vocabulaire \omterm{def:b2:affine:subspace}{affine}.
\end{remark}

\begin{example}[Géométrie barycentrique classique]\label{ex:b2:affine:median}
Le centre de gravité d'un triangle $ABC$ est le \omterm{def:b2:affine:barycenter}{barycentre} $G =
\operatorname{bar}(A,1; B,1; C,1)$. L'associativité avec le
milieu $A' = \operatorname{bar}(B, 1; C, 1)$ montre
\[
G = \operatorname{bar}(A, 1;\ A', 2) :
\]
$G$ est sur la médiane $AA'$ aux deux tiers de celle-ci --- et de même
pour les deux autres médianes~: les trois médianes sont concourantes, en
une ligne de calcul barycentrique.
\end{example}

\begin{example}[Les bimédianes d'un
quadrilatère]\label{ex:b2:affine:bimedians}
Soit $ABCD$ un quadrilatère quelconque (plan ou non~!) et considérons
ses \emph{bimédianes}~: les segments joignant les milieux des
côtés opposés, $M_{AB}M_{CD}$ et $M_{BC}M_{DA}$. Introduisons
le \omterm{def:b2:affine:barycenter}{barycentre} $G$ de $(A,1; B,1; C,1; D,1)$ et regroupons les
poids de deux façons~:
\[
G = \operatorname{bar}\bigl(M_{AB}, 2;\ M_{CD}, 2\bigr)
= \operatorname{bar}\bigl(M_{BC}, 2;\ M_{DA}, 2\bigr) :
\]
$G$ est le milieu des \emph{deux} bimédianes --- de sorte que les deux
bimédianes se coupent toujours en leur milieu, et le quadrilatère des
quatre milieux est un parallélogramme (ses diagonales sont les
bimédianes). Pas d'analyse de cas, pas de coordonnées, et
l'argument survit inchangé pour un quadrilatère gauche de
$\R^3$, où une démonstration fondée sur une figure serait déjà
délicate~: l'associativité se moque de la dimension.
\end{example}

\begin{example}[Classifier une application affine, de bout en
bout]\label{ex:b2:affine:dilation}
Soit $f(x, y) = (2x - 1,\ 3y - 4)$ sur $\R^2$. Sa partie linéaire
est $\varphi = \operatorname{diag}(2, 3)$, dont le \omterm{def:b2:reduction:eigen}{spectre}
$\{2, 3\}$ évite $1$~: par le critère de point fixe démontré
plus bas (\cref{prop:b2:affine:fixedpoint}), $f$ a exactement un
point fixe, obtenu en résolvant
\[
x = 2x - 1, \qquad y = 3y - 4
\qquad\Longrightarrow\qquad \Omega = (1, 2).
\]
En recentrant en $\Omega$ (poser $x = 1 + u$, $y = 2 + v$)~:
\[
f(1 + u,\ 2 + v) = (1 + 2u,\ 2 + 3v) :
\]
dans le repère en $\Omega$, $f$ \emph{est} sa partie linéaire, une
dilatation anisotrope d'un facteur $2$ horizontalement et $3$
verticalement à partir du centre $(1, 2)$. La leçon générale~: une
\omterm{def:b2:affine:subspace}{application affine} est «~application linéaire plus donnée de
position~», et la donnée de position se réduit à une origine bien
choisie dès que $1$ n'est pas \omterm{def:b2:reduction:eigen}{valeur propre}. Réciproquement,
translater mal l'origine \emph{crée} les termes constants~: la
géométrie \omterm{def:b2:affine:subspace}{affine} est l'art de choisir où placer $0$.
\end{example}

\begin{remark}[Méthode~: concours et alignement par
barycentres]
L'\cref{ex:b2:affine:median} est un cas particulier d'une recette
générale. Pour prouver que trois céviennes d'un triangle sont
concourantes, exhibons un \emph{unique} système pondéré $(A,
\alpha; B, \beta; C, \gamma)$ et utilisons l'associativité de trois
façons~: regrouper $(B, C)$ montre que le \omterm{def:b2:affine:barycenter}{barycentre} est sur la
cévienne issue de $A$, regrouper $(C, A)$ sur la cévienne issue de $B$,
regrouper $(A, B)$ sur la troisième. Pour les médianes, le système
$(A, 1; B, 1; C, 1)$ fait tout le travail~; pour des céviennes coupant
les côtés dans des rapports prescrits, les poids se lisent sur les
rapports. Pour prouver que trois points sont \emph{alignés}, en écrire
un comme \omterm{def:b2:affine:barycenter}{barycentre} des deux autres (\cref{exo:b2:affine:2}), ou
utiliser le critère du \omterm{def:b2:linalg:det}{déterminant} de l'\cref{exo:b2:affine:11}. Les
deux recettes remplacent l'ingéniosité géométrique par une comptabilité
de poids --- c'est précisément à cela que sert le calcul barycentrique.
\end{remark}

\section{La convexité, du point de vue affine}

\begin{definition}\label{def:b2:affine:convex}
Une partie $C$ d'un \omterm{def:b2:affine:def}{espace affine} est \emph{convexe}\index{ensemble convexe} lorsqu'elle contient
tout \omterm{def:b2:affine:barycenter}{barycentre} à poids \emph{positifs} de ses points ---
de manière équivalente, tout segment $\intcc{A}{B} = \{\operatorname{bar}(A,
1-t; B, t) : t \in \intcc{0}{1}\}$ entre ses points. L'\emph{enveloppe
convexe}\index{enveloppe convexe} $\operatorname{conv}(S)$ est l'ensemble de tous les
\omterm{def:b2:affine:barycenter}{barycentres} à poids positifs de points de $S$ --- la plus petite
partie convexe contenant $S$.
\end{definition}

\begin{example}[Les épigraphes sont des ensembles
convexes]\label{ex:b2:affine:epigraph}
La région $C = \{(x, y) : y \geq x^2\}$ au-dessus de la parabole
est \omterm{def:b2:affine:convex}{convexe}~: pour $(x_1, y_1), (x_2, y_2) \in C$ et $t \in
\intcc01$, l'inégalité de convexité de la fonction carré
donne
\[
\bigl((1-t)x_1 + tx_2\bigr)^2 \leq (1-t)x_1^2 + tx_2^2
\leq (1-t)y_1 + ty_2 ,
\]
donc le \omterm{def:b2:affine:barycenter}{barycentre} reste au-dessus de la parabole. Le calcul
est général~: $\{y \geq f(x)\}$ est \omterm{def:b2:affine:convex}{convexe} exactement lorsque $f$ est
une fonction \omterm{def:b2:affine:convex}{convexe} --- les \emph{ensembles} \omterm{def:b2:affine:convex}{convexes} et les
\emph{fonctions} \omterm{def:b2:affine:convex}{convexes} (\cref{ch:b2:realfun}) sont deux faces d'une
même notion, les épigraphes en étant le dictionnaire. C'est la raison
géométrique de l'existence des droites d'appui pour les fonctions
\omterm{def:b2:affine:convex}{convexes}, le fait qui démontrera l'inégalité de Jensen au
\cref{ch:b2:randomvar}.
\end{example}

\begin{example}[Générateurs redondants d'un
convexe]\label{ex:b2:affine:squarehull}
Soit $S = \{(0,0), (2,0), (2,2), (0,2), (1,1)\}$. Le cinquième
point est le \omterm{def:b2:affine:barycenter}{barycentre}
\[
(1,1) = \operatorname{bar}\bigl((0,0), \tfrac12;\ (2,2),
\tfrac12\bigr),
\]
il est donc déjà dans l'enveloppe des quatre autres~:
$\operatorname{conv}(S)$ est le carré ayant les quatre coins
pour sommets. En général, un point de $S$ qui est un
\omterm{def:b2:affine:barycenter}{barycentre} à poids positifs des \emph{autres} points de
$S$ peut être supprimé sans changer l'enveloppe~; les points qui
ne peuvent jamais être supprimés (ici les quatre coins) sont les
\emph{points extrémaux} de l'enveloppe. Les déterminer est un pur
calcul de \omterm{def:b2:affine:barycenter}{barycentre}~: $(2,0)$, par exemple, ne peut s'écrire comme
$\operatorname{bar}$ des points restants avec des poids positifs, car
la première coordonnée forcerait tout le poids sur les points avec $x =
2$, et la seconde coordonnée échoue alors. Les questions de convexité
se ramènent, encore et encore, à résoudre de petits systèmes pondérés.
\end{example}

\begin{theorem}[Carathéodory]\label{thm:b2:affine:caratheodory}
Dans un \omterm{def:b2:affine:def}{espace affine} de dimension $n$, tout point de
$\operatorname{conv}(S)$ est un \omterm{def:b2:affine:barycenter}{barycentre} d'au plus $n + 1$ points
de $S$.\index{théorème de Carathéodory}
\end{theorem}

\begin{proof}
Soit $G = \operatorname{bar}(A_0, \lambda_0; \dots; A_k, \lambda_k)$
avec $\lambda_i > 0$, $\sum\lambda_i = 1$, et $k + 1 > n + 1$
points. Les $k$ vecteurs $\vect{A_0A_i}$ ($i \geq 1$) sont liés
($k > n$)~: $\sum_{i\geq1}\mu_i \vect{A_0A_i} = 0$ de façon non
triviale~; en posant $\mu_0 = -\sum_{i\geq1}\mu_i$, on obtient des
poids $(\mu_i)$ avec $\sum\mu_i = 0$, $\sum \mu_i\,\vect{OA_i} = 0$
(pour tout $O$), non tous nuls. Alors pour tout réel $t$ les poids
$\lambda_i - t\mu_i$ sont encore de somme $1$ et, puisque
$\sum_i\mu_i\vect{OA_i} = 0$,
\[
\sum_i(\lambda_i - t\mu_i)\,\vect{OA_i}
= \sum_i\lambda_i\,\vect{OA_i} :
\]
ils produisent le \emph{même} point $G$. Faisons maintenant glisser
$t$ à partir de $0$~: un certain $\mu_i$ est positif (ils sont de
somme nulle et non tous nuls), donc
\[
t^* = \min\Bigl\{\frac{\lambda_i}{\mu_i} : \mu_i > 0\Bigr\}
\]
est bien défini et positif. En $t = t^*$~: pour les indices avec
$\mu_i > 0$, $\lambda_i - t^*\mu_i \geq 0$ par minimalité,
avec égalité en un indice minimisant~; pour les indices avec $\mu_i
\leq 0$, $\lambda_i - t^*\mu_i \geq \lambda_i > 0$. Tous
les poids restent positifs et au moins un s'est annulé~: $G$ se
réécrit comme \omterm{def:b2:affine:barycenter}{barycentre} de moins de points. On itère tant qu'il
reste plus de $n + 1$ points.
\end{proof}

\begin{example}
Dans le plan ($n = 2$)~: tout point de l'\omterm{def:b2:affine:convex}{enveloppe convexe} d'un ensemble
fini est dans un triangle dont les sommets sont dans l'ensemble --- le
contenu géométrique de Carathéodory, utilisé aussi bien en
optimisation qu'en probabilités (mélanges).
\end{example}

\begin{example}[Exécution de l'algorithme de
Carathéodory]\label{ex:b2:affine:caratheodoryrun}
Écrivons le centre du carré de
l'\cref{ex:b2:affine:squarehull} avec ses quatre coins $A_1 =
(0,0)$, $A_2 = (2,0)$, $A_3 = (2,2)$, $A_4 = (0,2)$~:
\[
(1,1) = \operatorname{bar}\bigl(A_1, \tfrac14;\ A_2,
\tfrac14;\ A_3, \tfrac14;\ A_4, \tfrac14\bigr),
\]
quatre points en dimension $2$ --- un de trop. La recette de la
démonstration demande des poids $(\mu_i)$ avec $\sum\mu_i = 0$ et
$\sum\mu_i\vect{OA_i} = 0$~: ici $\mu = (1, -1, 1, -1)$ convient
(les deux diagonales ont même milieu). Faire glisser $\lambda_i
\mapsto \lambda_i - t\mu_i$ laisse le \omterm{def:b2:affine:barycenter}{barycentre} fixe pour
tout $t$~; la valeur admissible extrémale $t = \frac14$ rend
les poids $(0, \tfrac12, 0, \tfrac12)$, annulant $A_1$ et
$A_3$ simultanément~:
\[
(1,1) = \operatorname{bar}\bigl(A_2, \tfrac12;\ A_4,
\tfrac12\bigr),
\]
une représentation par deux points --- encore mieux que les trois
que le théorème garantit, car le centre se trouve être
sur un segment entre générateurs. L'algorithme est
entièrement mécanique~: trouver une dépendance, faire glisser
jusqu'à ce qu'un poids s'annule, recommencer.
\end{example}

\section{Outils de classification affine}

\begin{proposition}[Points fixes des applications affines]\label{prop:b2:affine:fixedpoint}
Soit $f$ un endomorphisme \omterm{def:b2:affine:subspace}{affine} d'un \omterm{def:b2:affine:def}{espace affine} de dimension
finie, de partie linéaire $\varphi$. Si $1 \notin
\operatorname{Sp}(\varphi)$, alors $f$ a exactement un point fixe
$\Omega$, et dans la vectorialisation en $\Omega$, $f$ \emph{est}
sa partie linéaire. (Les translations, avec $\varphi = \mathrm{id}$ et
sans point fixe, sont l'obstruction fondamentale.)
\end{proposition}

\begin{proof}
Fixons $O$ et écrivons $f(O + x) = f(O) + \varphi(x)$. Le point
$O + x$ est fixe si et seulement si $O + x = f(O) + \varphi(x)$,
c'est-à-dire
\[
(\mathrm{id} - \varphi)(x) = \vect{O f(O)} .
\]
En dimension finie, $\mathrm{id} - \varphi$ est inversible
si et seulement si $0$ n'est pas \omterm{def:b2:reduction:eigen}{valeur propre} de $\mathrm{id} -
\varphi$, si et seulement si $1 \notin \operatorname{Sp}\varphi$ ---
et dans ce cas l'équation affichée a exactement une solution $x^*$,
donnant l'unique point fixe $\Omega = O + x^*$. Recentrons~: pour tout
vecteur $u$,
\[
f(\Omega + u) = f(\Omega) + \varphi(u) = \Omega + \varphi(u),
\]
donc dans le repère d'origine $\Omega$ l'application s'écrit $u \mapsto
\varphi(u)$~: purement linéaire. Lorsque $1 \in
\operatorname{Sp}\varphi$, soit aucun point fixe n'existe (l'équation
affichée peut être insoluble, comme pour une translation), soit tout
un sous-espace \omterm{def:b2:affine:subspace}{affine} de points fixes existe (ajouter à une solution
n'importe quel vecteur propre de \omterm{def:b2:reduction:eigen}{valeur propre} $1$)~: l'unicité est
exactement la condition spectrale.
\end{proof}

\begin{example}[Isométries du plan, complétées]\label{ex:b2:affine:isometries}
Une isométrie \omterm{def:b2:affine:subspace}{affine} du plan euclidien a sa partie linéaire dans
$O(2)$~: une rotation $R_\theta$ ou une réflexion (volume de l'année 1).
Si $\theta \neq 0$~: $1 \notin \operatorname{Sp} R_\theta$, donc
l'application est une \emph{rotation autour d'un unique centre}
(\cref{prop:b2:affine:fixedpoint}). Si la partie linéaire est une
réflexion~: soit une réflexion d'axe (des points fixes existent), soit
une \emph{réflexion glissée} (réflexion composée avec une translation
le long de l'axe, sans point fixe). Avec les translations, c'est la
classification complète des isométries du plan.
\end{example}

\begin{remark}[Les isométries du plan, en un coup d'œil]
Rassemblons les cas~: l'identité~; les translations ($\vec f =
\mathrm{id}$, sans point fixe sauf trivialité)~; les rotations
(partie linéaire $R_\theta$, $\theta \neq 0$~: un centre)~;
les réflexions (partie linéaire une réflexion, une droite de points
fixes)~; les réflexions glissées (même partie linéaire, sans point
fixe). Quatre familles plus l'identité, chacune reconnue par
deux données seulement~: la partie linéaire et l'ensemble des points
fixes --- le schéma de la \cref{prop:b2:affine:fixedpoint} rendu
exhaustif.
\end{remark}

\begin{example}[Une réflexion glissée, prise sur le
fait]\label{ex:b2:affine:glide}
Soit $f(x, y) = (y + 1,\ x + 1)$. La partie linéaire $(x, y)
\mapsto (y, x)$ est la réflexion d'axe la diagonale $y = x$, donc
$1 \in \operatorname{Sp}\vec f$ et
la \cref{prop:b2:affine:fixedpoint} reste muette. Des points fixes
exigeraient $x = y + 1$ et $y = x + 1$ simultanément~:
impossible --- il n'y en a aucun, donc $f$ n'est pas une réflexion.
Élever au carré règle la classification~:
\[
f\bigl(f(x, y)\bigr) = f(y + 1,\ x + 1) = (x + 2,\ y + 2),
\]
la translation de vecteur $(2, 2)$~: $f$ est la \emph{réflexion
glissée} d'axe la droite $y = x$ (décalée
convenablement~: le milieu de $M$ et $f(M)$ est toujours sur
$y = x + {}$constante, ici $y = x$, comme on le vérifie sur $M = (0,
0) \mapsto (1,1)$) et de vecteur de glissement $(1, 1)$, moitié de $f
\circ f$. À comparer avec l'\cref{exo:b2:affine:6}, où la même
partie linéaire mais une constante différente produisait une véritable
réflexion~: en présence de la \omterm{def:b2:reduction:eigen}{valeur propre} $1$, le terme constant
décide de tout.
\end{example}

\begin{example}[Les récurrences affines sont des dynamiques
affines]\label{ex:b2:affine:recursion}
La récurrence classique $u_{n+1} = au_n + b$ ($a \neq 1$)
itère l'\omterm{def:b2:affine:subspace}{application affine} $f(x) = ax + b$ de la droite, dont la
partie linéaire $a$ évite la \omterm{def:b2:reduction:eigen}{valeur propre} $1$~: il y a un unique
point fixe $\omega = \frac{b}{1-a}$, et y recentrer
(le cas unidimensionnel de la proposition ci-dessus) transforme
$f$ en la multiplication par $a$~:
\[
u_{n+1} - \omega = a\,(u_n - \omega)
\qquad\Longrightarrow\qquad
u_n = \omega + a^n(u_0 - \omega) .
\]
Pour $u_{n+1} = \frac{u_n}2 + 3$~: $\omega = 6$ et $u_n = 6 +
(u_0 - 6)2^{-n} \to 6$. La recette enseignée pour de telles
récurrences au \cref{ch:b2:series} --- «~soustraire le point fixe~» ---
est exactement la vectorialisation d'une \omterm{def:b2:affine:subspace}{application affine} en son
point fixe~; la convergence pour $\abs a < 1$ est le phénomène de
contraction que le \cref{ch:b2:metric} a transformé en théorème du
point fixe de Banach. Une idée, trois chapitres.
\end{example}

\begin{example}[Trouver le centre d'une rotation]\label{ex:b2:affine:rotationcenter}
Soit $f(x, y) = (-y + 2,\ x)$. La partie linéaire est
$\varphi(x, y) = (-y, x)$~: la rotation d'angle $\frac\pi2$,
dont le \omterm{def:b2:reduction:eigen}{spectre} $\{\iu, -\iu\}$ évite $1$. Par
la \cref{prop:b2:affine:fixedpoint} il y a exactement un point fixe~:
$x = -y + 2$ et $y = x$ donnent $x = 1$, $y = 1$, donc
$\Omega = (1, 1)$, et $f$ \emph{est} la rotation de centre
$(1, 1)$ et d'angle $\frac\pi2$. La leçon générale~: lorsque $1
\notin \operatorname{Sp}\vec f$, classifier $f$ coûte un
système linéaire --- la géométrie est entièrement dans la partie
linéaire, l'arithmétique entièrement dans la localisation du centre.
\end{example}

\begin{remark}[Où le langage affine sert ensuite]
Les \omterm{def:b2:affine:barycenter}{barycentres} et les \omterm{def:b2:affine:subspace}{applications affines} sont la grammaire des
chapitres de géométrie à venir~: les tangentes et les plans tangents
sont des objets \omterm{def:b2:affine:subspace}{affines}
(\cref{ch:b2:curves,ch:b2:surfaces}), un changement de
variables \omterm{def:b2:affine:subspace}{affine} multiplie les aires et les volumes par
$\abs{\det \vec f}$ (\cref{ch:b2:multint}), et l'espérance
est un \omterm{def:b2:affine:barycenter}{barycentre} dont les poids sont donnés par une loi de
probabilité, ce pourquoi la convexité gouverne l'inégalité de Jensen
(\cref{ch:b2:randomvar}). Dans le volume de l'année 3, le même
vocabulaire de convexité porte l'étude des \omterm{def:b2:nvs:norm}{normes} $L^p$ et des
inégalités intégrales.
\end{remark}

\begin{remark}[Perspectives dans ce volume]
Deux fils quittent ce chapitre. Le fil \emph{\omterm{def:b2:affine:subspace}{affine}}~:
les tangentes (\cref{ch:b2:curves}) et les plans tangents
(\cref{ch:b2:surfaces}) sont des sous-espaces \omterm{def:b2:affine:subspace}{affines} attachés à
des objets non linéaires, et la classification des quadriques dans le
chapitre sur les surfaces repose sur l'équation du centre
$A\Omega = -b$ de ce chapitre. Le fil \emph{\omterm{def:b2:affine:convex}{convexe}} est plus long~:
la convexité des demi-plans et des disques alimente la théorie de
Helly du problème du week-end~; la convexité des fonctions donne
l'inégalité de Jensen (\cref{ch:b2:randomvar})~; et le théorème final
du livre --- le critère d'extinction des processus de branchement
(\cref{ch:b2:genfun}) --- se décide par la position d'une
courbe \omterm{def:b2:affine:convex}{convexe} relativement à la diagonale, une image qui appartient
autant à ce chapitre qu'aux probabilités. Les \omterm{def:b2:affine:barycenter}{barycentres}
y reviennent aussi~: une espérance est un \omterm{def:b2:affine:barycenter}{barycentre} à poids de
probabilité.
\end{remark}

\section{Exercices}

\begin{exercise}[$\star$]\label{exo:b2:affine:1}
Dans $\R^3$, les ensembles suivants sont-ils des sous-espaces \omterm{def:b2:affine:subspace}{affines}~?
Donner directions et dimensions. $\{x + y + z = 1\}$~; $\;\{x + y + z =
1,\ x - z = 3\}$~; $\;\{x^2 + y^2 = 1\}$~; l'ensemble solution de $MX =
B$ pour un système compatible donné.
\end{exercise}

\begin{exercise}[$\star$]\label{exo:b2:affine:2}
Montrer que trois points distincts $A, B, C$ d'un \omterm{def:b2:affine:def}{espace affine} sont
alignés si et seulement si $C$ est un \omterm{def:b2:affine:barycenter}{barycentre} de $A$ et $B$, si et
seulement si les vecteurs $\vect{AB}, \vect{AC}$ sont liés. En déduire
une comptabilité de poids à la Ménélaüs~: si $C =
\operatorname{bar}(A, 1 - t; B, t)$, situer
$C$ pour $t = \frac12$, $t = 2$, $t = -1$.
\end{exercise}

\begin{exercise}[$\star$]\label{exo:b2:affine:3}
Soit $f$ l'\omterm{def:b2:affine:subspace}{application affine} de $\R^2$ donnée par $f(X) = MX + C$ avec
$M = \frac12\begin{pmatrix} 1 & 1\\ 1 & 1\end{pmatrix}$ et $C =
(1, 0)^{\mathsf T}$. Déterminer l'image de $f$, ses points fixes
(s'il y en a), et $f \circ f$.
\end{exercise}

\begin{exercise}[$\star\star$]\label{exo:b2:affine:4}
(L'associativité à l'œuvre) Dans un triangle $ABC$, soit $I, J, K$
divisant $BC$, $CA$, $AB$ dans les rapports $\vect{BI} =
\frac13\vect{BC}$, $\vect{CJ} = \frac13\vect{CA}$, $\vect{AK} =
\frac13\vect{AB}$. Exprimer $I, J, K$ comme \omterm{def:b2:affine:barycenter}{barycentres} et calculer le
\omterm{def:b2:affine:barycenter}{barycentre} de $(I,1;J,1;K,1)$~: que trouve-t-on, et pourquoi
était-ce prévisible~?
\end{exercise}

\begin{exercise}[$\star\star$]\label{exo:b2:affine:5}
Montrer qu'une application $f \colon \R^n \to \R^n$ conservant les
\emph{milieux} ($f\bigl(\frac{A+B}{2}\bigr) = \frac{f(A) + f(B)}{2}$)
et \emph{\omterm{def:b2:metric:continuity}{continue}} est \omterm{def:b2:affine:subspace}{affine}. \emph{(Montrer que l'application
vectorielle $u \mapsto f(O + u) - f(O)$ est additive via les milieux,
puis $\Q$-homogène, puis $\R$-homogène par continuité ---
la même stratégie de densité que pour l'équation fonctionnelle de
Cauchy dans le volume de l'année 1~; redémontrer ici les étapes
nécessaires.)}
\end{exercise}

\begin{exercise}[$\star\star$]\label{exo:b2:affine:6}
Classifier l'\omterm{def:b2:affine:subspace}{application affine} $f(x, y) = (y + 1,\; x - 1)$ du
plan euclidien~: partie linéaire, points fixes, nature géométrique
(réflexion~? réflexion glissée~?). Calculer $f \circ f$ et conclure.
\end{exercise}

\begin{exercise}[$\star\star\star$]\label{exo:b2:affine:7}
(Radon) Soit $A_1, \dots, A_{n+2}$ des points d'un \omterm{def:b2:affine:def}{espace affine} de
dimension $n$. Montrer qu'on peut les partager en deux groupes
disjoints dont les \omterm{def:b2:affine:convex}{enveloppes convexes} se rencontrent. \emph{(Comme
dans la démonstration de Carathéodory, trouver des poids $\mu_i$, non
tous nuls, avec $\sum\mu_i = 0$ et
$\sum\mu_i\vect{OA_i} = 0$~; séparer les poids positifs et négatifs et
normaliser les deux membres.)}
\end{exercise}

\begin{exercise}[$\star\star\star$]\label{exo:b2:affine:8}
Soit $f$ un endomorphisme \omterm{def:b2:affine:subspace}{affine} de $\R^n$ avec $f \circ f = f$.
Montrer que $f$ est la projection \omterm{def:b2:affine:subspace}{affine} sur le sous-espace \omterm{def:b2:affine:subspace}{affine}
$\operatorname{Fix}(f) = \operatorname{im} f$ parallèlement à la
direction $\ker\vec f$, et que réciproquement toutes ces projections
sont idempotentes. \emph{(Montrer d'abord que $\operatorname{im} f$ est
formé de points fixes.)}
\end{exercise}

\begin{exercise}[$\star$]\label{exo:b2:affine:9}
Soit $G = \operatorname{bar}(A, 1;\ B, 2;\ C, 3)$ dans un triangle
$ABC$. À l'aide de l'associativité, montrer que la droite $AG$ rencontre
$BC$ en $M = \operatorname{bar}(B, 2;\ C, 3)$, et situer $G$ sur le
segment $\intcc AM$~; situer de même l'intersection de $BG$
avec $CA$.
\end{exercise}

\begin{exercise}[$\star\star$]\label{exo:b2:affine:10}
Pour $\lambda \neq 0$, l'\emph{homothétie} $h_{\Omega,
\lambda}$ est l'\omterm{def:b2:affine:subspace}{application affine} fixant $\Omega$ de partie linéaire
$\lambda\,\mathrm{id}$. Montrer que la composée
$h_{\Omega', \mu} \circ h_{\Omega, \lambda}$ est une homothétie de
rapport $\lambda\mu$ lorsque $\lambda\mu \neq 1$, et une translation
lorsque $\lambda\mu = 1$~; dans le cas $\lambda = \mu = -1$ (deux
symétries centrales), calculer le vecteur de translation.
\end{exercise}

\begin{exercise}[$\star\star$]\label{exo:b2:affine:11}
(Ménélaüs) Dans un triangle $ABC$, soit $A' \in (BC)$, $B' \in
(CA)$, $C' \in (AB)$, tous distincts des sommets, et
définissons $\alpha, \beta, \gamma$ par $\vect{A'B} =
\alpha\,\vect{A'C}$, $\vect{B'C} = \beta\,\vect{B'A}$,
$\vect{C'A} = \gamma\,\vect{C'B}$. Montrer que $A', B', C'$ sont
alignés si et seulement si $\alpha\beta\gamma = 1$. \emph{(Écrire
chaque point comme \omterm{def:b2:affine:barycenter}{barycentre} de deux sommets~; montrer que trois
points sont alignés si et seulement si leurs lignes de coordonnées
barycentriques relativement à $(A, B, C)$ forment une matrice $3
\times 3$ singulière.)}
\end{exercise}

\begin{exercise}[$\star\star\star$]\label{exo:b2:affine:12}
Montrer que l'\omterm{def:b2:affine:convex}{enveloppe convexe} d'une partie compacte $K$ de $\R^n$
est compacte. \emph{(D'après le~\cref{thm:b2:affine:caratheodory},
$\operatorname{conv}(K)$ est l'image d'un \omterm{def:b2:metric:compact}{compact} par une
application \omterm{def:b2:metric:continuity}{continue}.)} Montrer sur un exemple dans $\R^2$ que
l'\omterm{def:b2:affine:convex}{enveloppe convexe} d'un ensemble \emph{fermé} n'est pas nécessairement
fermée.
\end{exercise}

\section{Problème~: de Radon à Helly, points centraux et théorème de
Jung}

\begin{omfigure}
\begin{tikzpicture}[scale=0.9]
  \draw[omDef, thick] (0,0) -- (3,0) -- (0.9,2.6) -- cycle;
  \fill (0,0) circle (1.6pt) node[below left] {\small $A_1$};
  \fill (3,0) circle (1.6pt) node[below right] {\small $A_2$};
  \fill (0.9,2.6) circle (1.6pt) node[above] {\small $A_3$};
  \fill[omThm] (1.3,0.9) circle (1.8pt)
    node[below, omThm] {\small $A_4$};
\end{tikzpicture}
\qquad
\begin{tikzpicture}[scale=0.9]
  \draw[omDef, thick] (0,0) -- (3,0.3) -- (3.3,2.4) --
    (0.4,2.2) -- cycle;
  \draw[omThm, thick] (0,0) -- (3.3,2.4);
  \draw[omThm, thick] (3,0.3) -- (0.4,2.2);
  \fill (0,0) circle (1.6pt) node[below left] {\small $A_1$};
  \fill (3,0.3) circle (1.6pt) node[below right] {\small $A_2$};
  \fill (3.3,2.4) circle (1.6pt) node[above right] {\small $A_3$};
  \fill (0.4,2.2) circle (1.6pt) node[above left] {\small $A_4$};
  \fill[omProp] (1.75,1.27) circle (1.8pt);
\end{tikzpicture}

{\small Les deux types de Radon pour quatre points du plan en
position générale~: un point à l'intérieur du triangle des autres
(partition $\{A_4\} \mid \{A_1, A_2, A_3\}$), ou position
\omterm{def:b2:affine:convex}{convexe}, où le point de Radon (orange) est l'intersection
des deux diagonales.}
\end{omfigure}

\begin{problem}[{Problème du week-end --- le théorème de Helly et deux
de ses dividendes}]\label{pb:b2:affine:1}
Le lemme de Radon (\cref{exo:b2:affine:7}) dit que $n + 2$
points d'un \omterm{def:b2:affine:def}{espace affine} de dimension $n$ se partagent toujours en
deux groupes aux \omterm{def:b2:affine:convex}{enveloppes convexes} qui se rencontrent. Ce problème
transforme ce seul fait d'\omterm{def:b2:structures:algebra}{algèbre} linéaire en une chaîne de théorèmes
de géométrie combinatoire~: le théorème d'intersection de Helly, le
théorème du point central (une médiane bidimensionnelle) et le
théorème de recouvrement de Jung. Partout, le plan est $\R^2$ muni de
sa structure euclidienne usuelle, et $\det$ est le \omterm{def:b2:linalg:det}{déterminant} dans
la base canonique.\index{théorème de Helly}\index{lemme de
Radon}\index{théorème de Jung}\index{point central}

\medskip
\noindent\textbf{Partie I --- Coordonnées barycentriques.} Les points
$A_0, \dots, A_k$ sont \emph{affinement indépendants} lorsque les
vecteurs $\vect{A_0A_1}, \dots, \vect{A_0A_k}$ sont linéairement
indépendants.
\begin{enumerate}
  \item Montrer que l'indépendance \omterm{def:b2:affine:subspace}{affine} ne dépend pas du
        choix du point de base $A_0$, et qu'elle équivaut à~:
        dès que deux familles de poids, chacune
        de somme $1$, définissent le même \omterm{def:b2:affine:barycenter}{barycentre} de $(A_0,
        \dots, A_k)$, les poids coïncident.
  \item Soit $(A, B, C)$ affinement indépendants dans le plan.
        Montrer que tout point $M$ admet un unique triplet
        $(\alpha, \beta, \gamma)$ avec $\alpha + \beta +
        \gamma = 1$ et $M = \operatorname{bar}(A, \alpha; B,
        \beta; C, \gamma)$ --- ses \emph{coordonnées
        barycentriques}.
  \item Démontrer les formules déterminantales
        \[
        \alpha = \frac{\det(\vect{MB},
        \vect{MC})}{\det(\vect{AB}, \vect{AC})},
        \qquad
        \beta = \frac{\det(\vect{MC},
        \vect{MA})}{\det(\vect{AB}, \vect{AC})},
        \qquad
        \gamma = \frac{\det(\vect{MA},
        \vect{MB})}{\det(\vect{AB}, \vect{AC})} :
        \]
        les coordonnées barycentriques sont des rapports d'aires
        signées.
  \item Les droites $BC$, $CA$, $AB$ sont les droites de
        coordonnées $\{\alpha = 0\}$, $\{\beta = 0\}$, $\{\gamma = 0\}$.
        Montrer que $M$ est dans le triangle fermé
        $\operatorname{conv}\{A, B, C\}$ si et seulement si $\alpha,
        \beta, \gamma \geq 0$, et que les trois droites découpent le
        plan en exactement sept régions, classées par les signes de
        $(\alpha, \beta, \gamma)$ (la configuration de signes $(-, -,
        -)$ étant impossible).
  \item Soit $u \colon \R^2 \to \R$ une \omterm{def:b2:affine:subspace}{application affine} (une
        \emph{forme \omterm{def:b2:affine:subspace}{affine}}). Montrer que $u(M) = \alpha\,u(A) +
        \beta\,u(B) + \gamma\,u(C)$, que les lignes de niveau d'une
        forme \omterm{def:b2:affine:subspace}{affine} non constante sont des droites, que toute droite
        s'obtient ainsi, et que les demi-plans fermés $\{u
        \geq c\}$ sont \omterm{def:b2:affine:convex}{convexes}.
\end{enumerate}

\medskip
\noindent\textbf{Partie II --- Partitions de Radon, précisées.} Une
famille de $n + 2$ points de $\R^n$ est en \emph{position
générale} lorsque tout sous-ensemble de $n + 1$ d'entre eux est
affinement indépendant. Une \emph{dépendance \omterm{def:b2:affine:subspace}{affine}} de $(A_1, \dots,
A_{n+2})$ est une famille $(\mu_i)$ avec $\sum_i \mu_i = 0$ et $\sum_i
\mu_i\,\vect{OA_i} = 0$ pour une (donc toute) origine $O$.
\begin{enumerate}[resume]
  \item Calculer une dépendance \omterm{def:b2:affine:subspace}{affine} non nulle des quatre
        points $A_1 = (0,0)$, $A_2 = (3,0)$, $A_3 = (0,3)$,
        $A_4 = (1,1)$~; donner la partition de Radon et le point
        de Radon.
  \item Montrer que pour des points en position générale l'espace
        vectoriel des dépendances \omterm{def:b2:affine:subspace}{affines} est de dimension exactement
        $1$, et qu'une dépendance non nulle n'a \emph{aucun}
        coefficient nul.
  \item En déduire que la partition de Radon de $n + 2$ points en
        position générale est unique (à l'échange des deux
        blocs près), chaque bloc étant l'ensemble des indices où
        $\mu_i$ a un signe fixé.
  \item Pour quatre points du plan en position générale,
        montrer la dichotomie~: soit la partition est de type
        $(1, 3)$ --- un point intérieur au triangle
        des trois autres --- soit de type $(2, 2)$~: les quatre points
        sont en position \omterm{def:b2:affine:convex}{convexe} et les segments joignant les
        deux paires (les diagonales) se coupent, au point de Radon.
  \item Reprendre la question 6 pour le carré unité $(0,0)$,
        $(1,0)$, $(1,1)$, $(0,1)$~: dépendance, partition,
        point de Radon.
\end{enumerate}

\medskip
\noindent\textbf{Partie III --- Le théorème de Helly dans le plan.}
\begin{enumerate}[resume]
  \item Soit $C_1, C_2, C_3, C_4$ des parties \omterm{def:b2:affine:convex}{convexes} de $\R^2$,
        trois quelconques d'entre elles ayant un point commun. Choisir
        $x_i \in \bigcap_{j \neq i} C_j$ et appliquer le lemme de
        Radon à $x_1, \dots, x_4$~: montrer que le point de Radon
        appartient aux quatre ensembles. \emph{(Pour chaque $k$, le
        bloc ne contenant pas $x_k$ est formé de points de $C_k$.)}
  \item (Helly) Soit $C_1, \dots, C_m$ ($m \geq 3$) des parties
        \omterm{def:b2:affine:convex}{convexes} de $\R^2$, trois quelconques d'entre elles
        s'intersectant. Montrer $\bigcap_{i=1}^m C_i \neq \emptyset$,
        par récurrence sur $m$~: remplacer $C_{m-1}$ et $C_m$ par
        $C_{m-1} \cap C_m$ et vérifier l'hypothèse pour la nouvelle
        famille à l'aide de la question 11.
  \item Trois contre-exemples, un par hypothèse~: (a) les
        trois côtés fermés d'un triangle s'intersectent deux à deux
        mais n'ont aucun point commun ($3$ ne peut être abaissé à
        $2$)~; (b) les quatre ensembles $S_i = \{x_1, \dots, x_4\}
        \setminus \{x_i\}$, pour quatre points en position
        générale, vérifient l'hypothèse d'intersection triple mais pas
        la conclusion (la convexité importe)~; (c) les demi-plans
        fermés $H_k = \intco{k}{+\infty} \times
        \R$, $k \in \N$, s'intersectent deux à deux et trois à trois
        mais $\bigcap_k H_k = \emptyset$ (les familles infinies
        exigent la compacité).
  \item (Helly \omterm{def:b2:metric:compact}{compact}) Soit $(K_i)_{i \in I}$ une famille
        quelconque de parties \omterm{def:b2:metric:compact}{compactes} \omterm{def:b2:affine:convex}{convexes} de $\R^2$, trois
        quelconques d'entre elles s'intersectant. À l'aide de la
        question 12 et de la propriété de Borel--Lebesgue
        (\cref{thm:b2:metric:borellebesgue}), montrer
        $\bigcap_{i \in I} K_i \neq \emptyset$.
  \item (Premier dividende) Soit $S$ un ensemble fini de points du
        plan et $r > 0$. Montrer~: si trois points quelconques de
        $S$ sont dans un disque fermé de rayon $r$, alors $S$
        est dans un disque fermé de rayon $r$. \emph{(Appliquer
        Helly aux disques $\overline D(p, r)$, $p \in S$.)}
\end{enumerate}

\medskip
\noindent\textbf{Partie IV --- Le théorème du point central.} Un
\emph{point central} d'un ensemble fini $S$ de $n$ points du
plan est un point $c$ (pas nécessairement dans $S$) tel que tout
demi-plan fermé contenant $c$ contient au moins $n/3$
points de $S$.
\begin{enumerate}[resume]
  \item (Dimension $1$) Pour des réels $x_1 \leq \dots \leq x_n$,
        montrer que la médiane $c = x_{\lceil n/2 \rceil}$
        vérifie~: toute demi-droite fermée contenant $c$
        contient au moins $n/2$ des $x_i$.
  \item (Lemme de dénombrement) Si $A, B, C$ sont des parties de $S$
        avec $\abs A, \abs B, \abs C > \tfrac{2n}3$, montrer $A \cap
        B \cap C \neq \emptyset$.
  \item Soit $m = \floor{2n/3} + 1$ et soit $\mathcal F$
        la famille (finie) des \omterm{def:b2:affine:convex}{enveloppes convexes}
        $\operatorname{conv}(T)$, $T \subseteq S$, $\abs T =
        m$. Montrer que trois membres quelconques de $\mathcal F$ ont
        un point commun, et déduire de Helly un point $c$
        commun à tous.
  \item Montrer que ce $c$ est un point central de $S$~: le
        \emph{théorème du point central}. \emph{(Si un demi-plan fermé
        passant par $c$ contenait moins de $n/3$
        points, son complémentaire \omterm{def:b2:metric:topology}{ouvert} contiendrait un ensemble $T$
        de $m$ points, et $\operatorname{conv}(T)$
        éviterait $c$.)}
  \item Optimalité~: soit $n = 3k$ et plaçons $k$ points dans
        chacun de trois disques de petit rayon $\varepsilon$ centrés
        aux sommets d'un grand triangle. Montrer que pour
        \emph{tout} point $c$ du plan un certain demi-plan
        fermé contenant $c$ contient au plus $n/3$
        points de $S$, de sorte que la constante $1/3$ ne peut être
        améliorée. \emph{(Parmi les trois directions de $c$
        vers les centres des disques, deux font un angle au plus
        $2\pi/3$.)}
\end{enumerate}

\medskip
\noindent\textbf{Partie V --- Le théorème de Jung et synthèse.}
\begin{enumerate}[resume]
  \item (Lemme du triangle) Soit $P, Q, R$ trois points de
        distances mutuelles $\leq 1$. Montrer qu'ils sont dans un
        disque fermé de rayon $1/\sqrt3$. \emph{(Si un
        angle est $\geq \pi/2$, prendre le disque de diamètre le
        plus long côté, à l'aide de la formule de la médiane
        $\norm{RM}^2 = \tfrac12\norm{RP}^2 +
        \tfrac12\norm{RQ}^2 - \tfrac14\norm{PQ}^2$~; si
        le triangle est acutangle, majorer le rayon du cercle
        circonscrit $a/(2\sin
        \widehat A)$ à l'aide de son plus grand angle, qui est dans
        $\intco{\pi/3}{\pi/2}$.)}
  \item (Jung) En déduire~: toute partie compacte du plan de
        diamètre $\leq 1$ est contenue dans un disque fermé de
        rayon $1/\sqrt3$.
  \item Optimalité~: pour le triangle équilatéral $A_1A_2A_3$
        de côté $1$ de centre de gravité $G$, démontrer l'identité
        de Leibniz $\sum_i \norm{\vect{OA_i}}^2 =
        3\norm{\vect{OG}}^2 + \sum_i \norm{\vect{GA_i}}^2$
        pour tout point $O$, et conclure que tout disque
        contenant les trois sommets a un rayon $\geq
        1/\sqrt3$, avec égalité seulement pour le disque circonscrit.
  \item (Helly dans $\R^n$) Énoncer et démontrer le théorème de Helly
        dans $\R^n$~: si un nombre fini d'\omterm{def:b2:affine:convex}{ensembles convexes} sont tels
        que $n + 1$ quelconques d'entre eux s'intersectent, alors tous
        s'intersectent. \emph{(Le lemme de Radon
        \cref{exo:b2:affine:7} traite $n + 2$ ensembles~; puis
        récurrence comme à la question 12.)}
  \item Synthèse. Assembler la chaîne
        \[
        \text{dépendance affine} \Rightarrow \text{Radon}
        \Rightarrow \text{Helly} \Rightarrow
        \text{point central et Jung},
        \]
        en indiquant en une phrase chacun~: où entre l'\omterm{def:b2:structures:algebra}{algèbre}
        linéaire, où entrent les signes des poids, où entre la
        convexité, et quelle unique étape a utilisé la dimension
        du plan. Que deviennent les constantes $3$ (dans Helly),
        $1/3$ (point central) et $1/\sqrt3$ (Jung) dans
        $\R^n$~? (Énoncer sans démonstration.)
\end{enumerate}
\end{problem}
